% GENERATED FILE. Edit canonical lessons, then run npm run generate:books.
% canonical-source-hash: fnv1a64:c5335359499d158e
% canonical-lessons: MR-C373-sarveksan, MR-C373-namuna, MR-C373-akdevari, MR-C373-vislesan, MR-C373-paddhat, MR-R373-review-366-373

\chapter{Data and Methods}
\label{ch:mr-data-and-methods}
\hlchaptermodality{373}

\begin{chapteropening}
\textbf{By the end of this chapter:} \emph{I can talk about surveys, samples, figures and methods.}

Complete the last lesson of chapter 373: I can talk about surveys, samples, figures and methods.
\end{chapteropening}

\section[\texorpdfstring{\mr{सर्वेक्षण} (sarvekṣaṇ)}{सर्वेक्षण (sarvekṣaṇ)}]{\mr{सर्वेक्षण} (sarvekṣaṇ) — a survey}

\label{lesson:MR-C373-sarveksan}



\begin{quote}
Before the new one: say the Marathi for a thesis, then the Marathi for a theory.
\end{quote}

\subsection*{You'll want to know: \mr{सर्वेक्षण}}
\textbf{\mr{सर्वेक्षण}} — \emph{sarvekṣaṇ} — \textquotedblleft{}a survey\textquotedblright{}.

\begin{quote}
\mr{संस्थेने} \mr{गावात} \mr{आरोग्याचं} \mr{सर्वेक्षण} \mr{केलं}.
\end{quote}

\begin{quote}
\emph{saṁsthene gāvāt ārogyācaṁ sarvekṣaṇ kelaṁ.}
\end{quote}

\begin{quote}
The organization carried out a health survey in the village.
\end{quote}

Neuter.

\mr{सर्व} + \mr{ईक्षण}, all + looking.

\begin{cousinweb}
\textbf{\mr{सर्वेक्षण}} (\emph{sarvekṣaṇ}) is a learned compound built from Sanskrit \emph{sarva}, \textquotedblleft{}all\textquotedblright{}, and \emph{īkṣaṇa}, \textquotedblleft{}look, view\textquotedblright{} (Monier-Williams): a view over everything.
\end{cousinweb}

\begin{grammarlens}[title={What you've built}]
One more word for giving reasons and opinions.
\end{grammarlens}

\subsection*{Your turn}
\begin{itemize}
\raggedright
  \item \emph{Say it:} \emph{sarvekṣaṇ}
  \item \emph{Say it:} \emph{sarvekṣaṇ}, once more
  \item \emph{Say it:} \emph{saṁsthene gāvāt ārogyācaṁ sarvekṣaṇ kelaṁ.} — The organization carried out a health survey in the village.
  \item \emph{Recall:} say the Marathi for a thesis, then the Marathi for a theory, then say \emph{sarvekṣaṇ} again
\end{itemize}

\begin{tcolorbox}[breakable,colback=teal!4,colframe=teal!35!black,title={Before you move on}]
What does \mr{सर्वेक्षण} mean? (\textquotedblleft{}A survey\textquotedblright{}.) Say it once more.
\end{tcolorbox}
\section[\texorpdfstring{\mr{नमुना} (namunā)}{नमुना (namunā)}]{\mr{नमुना} (namunā) — a sample; a specimen}

\label{lesson:MR-C373-namuna}



\begin{quote}
Before the new one: say the Marathi for a theory, then the Marathi for a survey.
\end{quote}

\subsection*{You'll want to know: \mr{नमुना}}
\textbf{\mr{नमुना}} — \emph{namunā} — \textquotedblleft{}a sample; a specimen\textquotedblright{}.

\begin{quote}
\mr{प्रयोगासाठी} \mr{नदीच्या} \mr{पाण्याचा} \mr{नमुना} \mr{घेतला}.
\end{quote}

\begin{quote}
\emph{prayogāsāṭhī nadīcyā pāṇyācā namunā ghetlā.}
\end{quote}

\begin{quote}
A sample of river water was taken for the experiment.
\end{quote}

Masculine. Plural: \mr{नमुने}, \emph{namune}.

Plural \mr{नमुने}; also a model, a pattern.

\begin{cousinweb}
\textbf{\mr{नमुना}} (\emph{namunā}) is a loan from Persian \emph{namūna}, \textquotedblleft{}a sample, a pattern\textquotedblright{}.
\end{cousinweb}

\begin{grammarlens}[title={What you've built}]
One more word for giving reasons and opinions.
\end{grammarlens}

\subsection*{Your turn}
\begin{itemize}
\raggedright
  \item \emph{Say it:} \emph{namunā}
  \item \emph{Say it:} \emph{namunā}, once more
  \item \emph{Say it:} \emph{prayogāsāṭhī nadīcyā pāṇyācā namunā ghetlā.} — A sample of river water was taken for the experiment.
  \item \emph{Recall:} say the Marathi for a theory, then the Marathi for a survey, then say \emph{namunā} again
\end{itemize}

\begin{tcolorbox}[breakable,colback=teal!4,colframe=teal!35!black,title={Before you move on}]
What does \mr{नमुना} mean? (\textquotedblleft{}A sample; a specimen\textquotedblright{}.) Say it once more.
\end{tcolorbox}
\section[\texorpdfstring{\mr{आकडेवारी} (ākḍevārī)}{आकडेवारी (ākḍevārī)}]{\mr{आकडेवारी} (ākḍevārī) — statistics; figures}

\label{lesson:MR-C373-akdevari}



\begin{quote}
Before the new one: say the Marathi for a survey, then the Marathi for a sample.
\end{quote}

\subsection*{You'll want to know: \mr{आकडेवारी}}
\textbf{\mr{आकडेवारी}} — \emph{ākḍevārī} — \textquotedblleft{}statistics; figures\textquotedblright{}.

\begin{quote}
\mr{सर्वेक्षणाची} \mr{आकडेवारी} \mr{उद्या} \mr{जाहीर} \mr{होईल}.
\end{quote}

\begin{quote}
\emph{sarvekṣaṇācī ākḍevārī udyā jāhīr hoīl.}
\end{quote}

\begin{quote}
The survey's figures will be announced tomorrow.
\end{quote}

Feminine.

\mr{आकडा}, a figure, + -\mr{वारी}: figures taken together.

\begin{cousinweb}
\textbf{\mr{आकडेवारी}} (\emph{ākḍevārī}) is built on \mr{आकडा}, \textquotedblleft{}a figure, a number\textquotedblright{}, with the ending -\mr{वारी}; the origin of \mr{आकडा} is not settled.
\end{cousinweb}

\begin{grammarlens}[title={What you've built}]
One more word for giving reasons and opinions.
\end{grammarlens}

\subsection*{Your turn}
\begin{itemize}
\raggedright
  \item \emph{Say it:} \emph{ākḍevārī}
  \item \emph{Say it:} \emph{ākḍevārī}, once more
  \item \emph{Say it:} \emph{sarvekṣaṇācī ākḍevārī udyā jāhīr hoīl.} — The survey's figures will be announced tomorrow.
  \item \emph{Recall:} say the Marathi for a survey, then the Marathi for a sample, then say \emph{ākḍevārī} again
\end{itemize}

\begin{tcolorbox}[breakable,colback=teal!4,colframe=teal!35!black,title={Before you move on}]
What does \mr{आकडेवारी} mean? (\textquotedblleft{}Statistics; figures\textquotedblright{}.) Say it once more.
\end{tcolorbox}
\section[\texorpdfstring{\mr{विश्लेषण} (viśleṣaṇ)}{विश्लेषण (viśleṣaṇ)}]{\mr{विश्लेषण} (viśleṣaṇ) — an analysis}

\label{lesson:MR-C373-vislesan}



\begin{quote}
Before the new one: say the Marathi for a sample, then the Marathi for statistics.
\end{quote}

\subsection*{You'll want to know: \mr{विश्लेषण}}
\textbf{\mr{विश्लेषण}} — \emph{viśleṣaṇ} — \textquotedblleft{}an analysis\textquotedblright{}.

\begin{quote}
\mr{तिने} \mr{आकडेवारीचं} \mr{नीट} \mr{विश्लेषण} \mr{केलं}.
\end{quote}

\begin{quote}
\emph{tine ākḍevārīcaṁ nīṭ viśleṣaṇ kelaṁ.}
\end{quote}

\begin{quote}
She analysed the figures carefully.
\end{quote}

Neuter.

\mr{विश्लेषण} \mr{करणे}, to analyse.

\begin{cousinweb}
\textbf{\mr{विश्लेषण}} (\emph{viśleṣaṇ}) is a tatsama, taken from Sanskrit \emph{viśleṣaṇa}, \textquotedblleft{}separating, dissolving\textquotedblright{} (Monier-Williams).
\end{cousinweb}

\begin{grammarlens}[title={What you've built}]
One more word for giving reasons and opinions.
\end{grammarlens}

\subsection*{Your turn}
\begin{itemize}
\raggedright
  \item \emph{Say it:} \emph{viśleṣaṇ}
  \item \emph{Say it:} \emph{viśleṣaṇ}, once more
  \item \emph{Say it:} \emph{tine ākḍevārīcaṁ nīṭ viśleṣaṇ kelaṁ.} — She analysed the figures carefully.
  \item \emph{Recall:} say the Marathi for a sample, then the Marathi for statistics, then say \emph{viśleṣaṇ} again
\end{itemize}

\begin{tcolorbox}[breakable,colback=teal!4,colframe=teal!35!black,title={Before you move on}]
What does \mr{विश्लेषण} mean? (\textquotedblleft{}An analysis\textquotedblright{}.) Say it once more.
\end{tcolorbox}
\section[\texorpdfstring{\mr{पद्धत} (paddhat)}{पद्धत (paddhat)}]{\mr{पद्धत} (paddhat) — a method; a way}

\label{lesson:MR-C373-paddhat}



\begin{quote}
Before the new one: say the Marathi for statistics, then the Marathi for an analysis.
\end{quote}

\subsection*{You'll want to know: \mr{पद्धत}}
\textbf{\mr{पद्धत}} — \emph{paddhat} — \textquotedblleft{}a method; a way\textquotedblright{}.

\begin{quote}
\mr{शिकवण्याची} \mr{त्यांची} \mr{पद्धत} \mr{वेगळी} \mr{आहे}.
\end{quote}

\begin{quote}
\emph{śikvaṇyācī tyāñcī paddhat vegḷī āhe.}
\end{quote}

\begin{quote}
Their way of teaching is different.
\end{quote}

Feminine.

Feminine: X \mr{ची} \mr{पद्धत}, the way of X.

\begin{cousinweb}
\textbf{\mr{पद्धत}} (\emph{paddhat}) is a tatsama, taken from Sanskrit \emph{paddhati}, \textquotedblleft{}a path, a track, a way\textquotedblright{} (Monier-Williams). Marathi drops the final -i.
\end{cousinweb}

\begin{grammarlens}[title={What you've built}]
One more word for giving reasons and opinions.
\end{grammarlens}

\subsection*{Your turn}
\begin{itemize}
\raggedright
  \item \emph{Say it:} \emph{paddhat}
  \item \emph{Say it:} \emph{paddhat}, once more
  \item \emph{Say it:} \emph{śikvaṇyācī tyāñcī paddhat vegḷī āhe.} — Their way of teaching is different.
  \item \emph{Recall:} say the Marathi for statistics, then the Marathi for an analysis, then say \emph{paddhat} again
\end{itemize}

\begin{tcolorbox}[breakable,colback=teal!4,colframe=teal!35!black,title={Before you move on}]
What does \mr{पद्धत} mean? (\textquotedblleft{}A method; a way\textquotedblright{}.) Say it once more.
\end{tcolorbox}
\section[(dialogue)]{Looking back: from Values to Data and Methods}

\label{lesson:MR-R373-review-366-373}



\begin{quote}
Say the Marathi for a value and morality.
\end{quote}

\subsection*{Your turn}
Say these aloud:

\begin{itemize}
\raggedright
  \item a value, morality, integrity, renunciation, welfare
  \item a practice, a convention, an ideology, a tendency, social
  \item to comply with, to violate, to discharge, to entrust, to handle
  \item a professor, a syllabus, a scholarship, a fee, a hostel
  \item a lecture, an essay, a speech, a symposium, a speaker
  \item to pass, to fail, to revise, to cram, a certificate
  \item research, an experiment, a scholar, a thesis, a theory
  \item a survey, a sample, statistics, an analysis, a method
\end{itemize}

\begin{tcolorbox}[breakable,colback=teal!4,colframe=teal!35!black,title={Before you move on}]
A value? (\textbf{\mr{मूल्य}}, \emph{mūlya}.) Morality? (\textbf{\mr{नैतिकता}}, \emph{naitiktā}.) Integrity? (\textbf{\mr{सचोटी}}, \emph{sacoṭī}.) Renunciation? (\textbf{\mr{त्याग}}, \emph{tyāg}.) Welfare? (\textbf{\mr{हित}}, \emph{hit}.) A practice? (\textbf{\mr{प्रथा}}, \emph{prathā}.) A convention? (\textbf{\mr{रूढी}}, \emph{rūḍhī}.) An ideology? (\textbf{\mr{विचारसरणी}}, \emph{vicārsarṇī}.) A tendency? (\textbf{\mr{प्रवृत्ती}}, \emph{pravṛttī}.) Social? (\textbf{\mr{सामाजिक}}, \emph{sāmājik}.) To comply with? (\textbf{\mr{पालन} \mr{करणे}}, \emph{pālan karṇe}.) To violate? (\textbf{\mr{उल्लंघन} \mr{करणे}}, \emph{ullaṅghan karṇe}.) To discharge? (\textbf{\mr{बजावणे}}, \emph{bajāvṇe}.) To entrust? (\textbf{\mr{सोपवणे}}, \emph{sopavṇe}.) To handle? (\textbf{\mr{हाताळणे}}, \emph{hātāḷṇe}.) A professor? (\textbf{\mr{प्राध्यापक}}, \emph{prādhyāpak}.) A syllabus? (\textbf{\mr{अभ्यासक्रम}}, \emph{abhyāsakram}.) A scholarship? (\textbf{\mr{शिष्यवृत्ती}}, \emph{śiṣyavṛttī}.) A fee? (\textbf{\mr{शुल्क}}, \emph{śulka}.) A hostel? (\textbf{\mr{वसतिगृह}}, \emph{vasatigṛh}.) A lecture? (\textbf{\mr{व्याख्यान}}, \emph{vyākhyān}.) An essay? (\textbf{\mr{निबंध}}, \emph{nibandh}.) A speech? (\textbf{\mr{भाषण}}, \emph{bhāṣaṇ}.) A symposium? (\textbf{\mr{परिसंवाद}}, \emph{parisaṁvād}.) A speaker? (\textbf{\mr{वक्ता}}, \emph{vaktā}.) To pass? (\textbf{\mr{उत्तीर्ण} \mr{होणे}}, \emph{uttīrṇa hoṇe}.) To fail? (\textbf{\mr{नापास} \mr{होणे}}, \emph{nāpās hoṇe}.) To revise? (\textbf{\mr{उजळणी} \mr{करणे}}, \emph{ujaḷṇī karṇe}.) To cram? (\textbf{\mr{घोकणे}}, \emph{ghokṇe}.) A certificate? (\textbf{\mr{प्रमाणपत्र}}, \emph{pramāṇpatra}.) Research? (\textbf{\mr{संशोधन}}, \emph{saṁśodhan}.) An experiment? (\textbf{\mr{प्रयोग}}, \emph{prayog}.) A scholar? (\textbf{\mr{विद्वान}}, \emph{vidvān}.) A thesis? (\textbf{\mr{प्रबंध}}, \emph{prabandh}.) A theory? (\textbf{\mr{सिद्धांत}}, \emph{siddhānt}.) A survey? (\textbf{\mr{सर्वेक्षण}}, \emph{sarvekṣaṇ}.) A sample? (\textbf{\mr{नमुना}}, \emph{namunā}.) Statistics? (\textbf{\mr{आकडेवारी}}, \emph{ākḍevārī}.) An analysis? (\textbf{\mr{विश्लेषण}}, \emph{viśleṣaṇ}.) A method? (\textbf{\mr{पद्धत}}, \emph{paddhat}.)
\end{tcolorbox}
